\chapter{Từ Lớp Học Ra Cuộc Đời}

\begin{leadbox}
Chương cuối không khép cuốn sách lại bằng những lời khuyên chung chung.
Nó mở rộng từ lớp học ra cuộc đời: những gì một học sinh học được về kỷ luật, tôn trọng, trung thực và lòng tốt rồi sẽ đi cùng em vào công việc, gia đình và cách đối xử với người khác.
\end{leadbox}

\storycard{111. Ra Trường Rồi Mới Hiểu}{Only Understood After Graduation}{Nhiều lời nhắc từng làm ta khó chịu, về sau lại trở thành điểm tựa giữa những năm tháng đi làm.}
\storycard{112. Gặp Lại Thầy Cũ}{Meeting the Old Teacher Again}{Khoảnh khắc gặp lại là lúc người ta nhìn rõ hơn mình đã từng được uốn nắn và cứu khỏi lệch hướng như thế nào.}
\storycard{113. Người Trò Từng Gây Rối Đi Dạy Lại}{The Troublemaker Who Becomes a Mentor}{Một quá khứ xấu không cấm ai trở thành người dẫn đường tốt nếu họ biết chuyển hóa nó thành kinh nghiệm sống.}
\storycard{114. Lời Cảm Ơn Chưa Từng Nói}{The Thank-You Never Said}{Có những món nợ tình cảm chỉ được nhận ra khi con người đã đủ trưởng thành để gọi đúng tên của nó.}
\storycard{115. Vết Thương Cần Được Sửa}{Wounds That Need Repair}{Không phải mọi ký ức học đường đều đẹp; trưởng thành cũng là học cách nhận diện và hàn lại những vết đau cũ.}
\storycard{116. Dạy Con Bằng Bài Học Đã Trả Giá}{Teaching the Next Generation}{Người lớn thường dạy trẻ bằng chính những bài học mà ngày xưa mình đã phải trả giá mới hiểu được.}
\storycard{117. Công Việc Cần Kỷ Luật Như Lớp Học}{Work Needs Discipline Too}{Những kỹ năng học đường tưởng nhỏ nhiều khi lại trở thành năng lực nghề nghiệp bền vững của một đời người.}
\storycard{118. Trung Thực Khi Không Ai Kiểm Tra}{Honesty Without Supervision}{Rời khỏi trường lớp, bài thi lớn nhất của nhân cách mới thật sự bắt đầu ở những nơi không có ai chấm điểm.}
\storycard{119. Tử Tế Với Người Yếu Hơn Mình}{Kindness Toward the Weaker}{Giá trị của giáo dục lộ rõ nhất trong cách con người sử dụng quyền lực với người ở thế yếu hơn mình.}
\storycard{120. Định Nghĩa Lại Hai Chữ Ngoan Và Tốt}{Redefining Good}{Khép lại cuốn sách là một cách hiểu sâu hơn: ngoan không phải phục tùng mù quáng, và tốt không có nghĩa là vô khuyết, mà là biết sửa mình và tôn trọng người khác.}
